\clearpage
\section*{Production Tables}

\begin{table}[htbp]
\centering
\small
\caption{Formal objects and tribunal outputs.}
\label{tab:formal-objects}
\begin{tabular}{p{0.20\linewidth}p{0.70\linewidth}}
\toprule
Object & Operational definition \\
\midrule
Jurisdiction $\mathcal J$ & Frozen semantic signature, admissible substitutions and interventions, contexts, tolerances and costs. \\
Carrier $C=(V,R,\Theta)$ & Components, typed/signed relations and the physical or dynamical regime required by the signature. \\
Closure & All relations required for the declared signature are counted within the candidate rather than supplied by an unregistered external mechanism. \\
Sufficiency & The complete strict signature is realized within tolerance. \\
Inclusion-minimality & No proper admissible subcarrier realizes the complete signature. \\
Invariance & The signature survives every admissible substitution and transport declared by the jurisdiction. \\
Positive & Closed, sufficient, inclusion-minimal and invariant carrier. \\
Negative & An observed hard failure of baseline sufficiency or a binding gate. \\
Abstain & Evidence is incomplete or does not identify one carrier; no hard promotion is made. \\
Class-valued & More than one non-equivalent minimal carrier survives. \\
\bottomrule
\end{tabular}
\end{table}

\begin{table}[htbp]
\centering
\scriptsize
\caption{Relation between the UMI tribunal and established frameworks. A check mark denotes a defining or usual component rather than an absolute property of every implementation.}
\label{tab:comparators}
\begin{tabular}{lccccccc}
\toprule
Framework & Topology & Dynamics & Ablation/rewire & Substitution & Transport & Minimality & Abstention \\
\midrule
Network motifs & $\checkmark$ & partial & no & no & no & no & no \\
Dynamical modules & partial & $\checkmark$ & partial & partial & partial & partial & no \\
Minimal realization & no & $\checkmark$ & no & equivalence & representation & $\checkmark$ & no \\
Bisimulation & no & $\checkmark$ & no & equivalence & representation & partial & no \\
Internal-model theory & typed & $\checkmark$ & necessity & plant variation & implementations & function-specific & no \\
Organisational closure & relational & $\checkmark$ & partial & partial & biological & not generally & no \\
UMI tribunal & candidate only & $\checkmark$ & $\checkmark$ & $\checkmark$ & $\checkmark$ & $\checkmark$ & $\checkmark$ \\
\bottomrule
\end{tabular}
\end{table}

\begin{table}[htbp]
\centering
\small
\caption{Evidence hierarchy retained in the cross-domain portfolio.}
\label{tab:evidence-hierarchy}
\begin{tabular}{lp{0.69\linewidth}}
\toprule
Code & Evidence type \\
\midrule
T1 & Exact theorem, algebraic identity or exhaustive finite result. \\
C1 & Reproducible computation with a frozen model and independent verifier. \\
M1 & Explicit mechanistic model whose assumptions and jurisdiction are declared. \\
P1 & Direct physical experiment observing the declared signature. \\
P2 & Physical perturbation, ablation or rewiring relevant to necessity. \\
P3 & Independent physical or computational substrate realizing the relation class. \\
S1 & Controlled literature synthesis linking non-identical implementations. \\
A & Ambiguous, incomplete or causally non-identifying evidence. \\
\bottomrule
\end{tabular}
\end{table}

\begin{table}[htbp]
\centering
\small
\caption{Claim boundary.}
\label{tab:claim-boundary}
\begin{tabular}{p{0.30\linewidth}p{0.27\linewidth}p{0.33\linewidth}}
\toprule
Established antecedent & Use in this work & Claim explicitly withheld \\
\midrule
Network motifs and functional modules & Candidate generation and comparison & First discovery of relational building blocks \\
Minimal realization and bisimulation & Equivalence/minimality comparators & Replacement of complete behavioral-equivalence theory \\
Internal-model and RPA module theory & High-grade domain evidence & Discovery of integral control or adaptation modules \\
RAF and autocatalytic-set theory & Exact chemical closure case & Discovery of autocatalytic closure \\
Topological bulk--boundary theory & Topological carrier case & Discovery of topological protection \\
Interventionist mechanism analysis & Necessity testing & First use of perturbation to identify mechanism \\
Cross-domain tribunal & Conjunctive synthesis and validation & Universal ontology, unique UMI for every system, or empirical multifractality \\
\bottomrule
\end{tabular}
\end{table}
